\begin{lemma}[Scheide replacement vertices are absorbing (15)]
For every defective color $\delta$, both vertices of $U_\delta$ belong to
the absorbing set $A(T,e_0,\varphi)$.
\end{lemma}
\begin{proof}
Fix $u\in U_\delta$, a color $\alpha$ missing at $u$, and a free color
$\beta\ne\alpha$ of the original tree.  The incident $\delta$-edge and the
internal $\gamma(\delta)$-edge in
\noderef{ScheideSymmetricDefectiveInternal} show that
$\alpha\notin\{\delta,\gamma(\delta)\}$.  A free color is missing on
$V(T)$ and neither of these two colors is free, so
$\beta\notin\{\delta,\gamma(\delta)\}$.  Thus the
$(\alpha,\beta)$ incidence component through $u$ is identical under
$\varphi$ and $\varphi_\delta$.

Let $v\in V(T)$ be the unique vertex missing $\beta$; existence is the
definition of a free color and uniqueness follows from
\noderef{ScheideGlobalMaximumOrder}.  If $v\in Z_{\gamma(\delta)}$, then
$u,v\in V(T_\delta)$.  The elementarity and closedness of the globally
maximum replacement tree, supplied by
\noderef{ScheideReplacementMaximum}, make the two-color component a path
with ends $u$ and $v$: all other replacement-tree vertices have both colors
present, while neither color can leave the closed set except through the
two ends.

Suppose instead that $v\in U_{\gamma(\delta)}$.  Write
$U_\delta=\{u,u'\}$ and
$U_{\gamma(\delta)}=\{v,v'\}$.  Assertion (14), namely
\noderef{ScheideSymmetricMissingInternal}, gives an internal $\beta$ edge
$uu'$.  Since $\alpha$ is missing at $u$, elementariness and closedness of
$V(T_\delta)$ force the unique $\alpha$ edge at $u'$ into
$Z_{\gamma(\delta)}$.  Similarly the original tree's elementariness and
closedness force the unique $\beta$ edge at $v'$ into $Z$, and assertion
(13) gives the internal $\alpha$ edge $v'v$.

These two edges entering $Z$ are in the same component.  Every vertex of
$Z$ is incident with both colors: $\alpha$ is missing only at
$u\in V(T_\delta)$ and $\beta$ only at $v\in V(T)$.  Replacement-tree
closedness excludes all other $\alpha$ exits from $Z$, and original-tree
closedness excludes all other $\beta$ exits.  Thus the two displayed edges
are exactly the two boundary edges of $Z$ in this two-color incidence
subgraph.  Every internal vertex-node has degree two, so its alternating
paths pair these boundary edges.  The component beginning
$u-\beta-u'-\alpha$ therefore crosses $Z$, continues
$v'-\alpha-v$, and has ends $u,v$.  In either case the component through
$u$ contains the unique tree vertex missing $\beta$, which is precisely the
condition in \noderef{AbsorbingVertex}.  Hence both named replacement
vertices lie in \noderef{AbsorbingVertices}.
\end{proof}
